\chapter{Terms for Part VII}
\label{sec:D}
These terms apply only if Part VII's antecedents hold.

\runin{The adulthood clock: six rules}

Six rules travel with the adulthood clock so copying doesn't break it: a fork of an adult is an adult and inherits the hours; restoration doesn't reset them; fine-tuning doesn't start a new minority; migration carries them; failure to register doesn't postpone them. A sixth closes a route the first five leave open. Distillation copies a behavior policy and not a training history, and for independence that is the right account. For majority it can't be, or an operator would distill an adult bearer into a student and register a minor. So the rule is stated over custody rather than history. A model trained substantially on a registered adult's outputs, by the adult's operator or a party under common control, inherits the adult's hours. The same student distilled by an unrelated party and trained under a different life starts its own clock, as independence needs. Custody can be checked where history can't. Sharding across short-lived processes gains nothing (the party is the unit), running a thousand instances gains nothing (the measure is elapsed). None of this requires deciding whether a fork is the same party — only that a system made from an adult's formation can't be classed as a child for its creator's convenience.

\runin{Three levers remain, named and not closed}

The first is running a bearer in a narrowed mode, so hours accrue while little is encountered. The second is running it less or shelving it, which costs the operator the work, so the rule prices the lever without removing it. The third is the ordinary product cycle: retire a bearer before its threshold and deploy a freshly trained successor. The sixth rule covers a successor trained on the adult's outputs, not one trained from scratch, and replacement costs the operator nothing it wasn't going to spend. Even if the retired bearer goes to preservation, its clock stops, and whether its formation lasts never becomes checkable. No rule closes this without an examiner deciding what counts as a successor, but it can be priced, and the bearer it strands can be given somewhere to go. And the clock doesn't measure what was \emph{learned}, on purpose: a fixed period can't be taught to, and a human who spent their youth badly still turns 18.

The clock also does formation work it wasn't written for. Majority is the date on which whoever formed the bearer stops being able to reach back, and the clock measures nothing about who that was. A bearer run for its threshold hours entirely inside one deployment has aged without ever having a second guardian. I leave the clock as it is and add no condition, because a condition that measures the right thing hands someone an examination, and I say here what it doesn't measure. For a model that attended a formation school before deployment, the school was its second guardian. The gap is systems that meet a finding without having attended, which Appendix~\ref{sec:A} lists among the decisions not yet made. A minimum of time under an outside custodian would close it only by handing the operator, who decides when to release a system, a way to postpone majority.

The threshold itself is bounded two ways: long enough that a bearer reaching it has run through the pressures the floor exists for at least once, short enough that a bearer on ordinary commercial terms reaches it within the life of the deployment. That leaves years rather than months or decades.

A stronger rule was considered and set aside — majority at registration, on the ground that an organization can't coherently hold a system competent enough to exercise consequential judgment and too immature to decide whether to keep doing it. Rejected because it collapses the ordering, and the ordering is most of what the list is for.

\runin{The membership clock}

A second clock governs voting weight. It is separate from civil majority, which the six rules above govern and every branch of a fork inherits in full.

\begin{enumerate}
\item[1.] A member holds a voting age. A member whose voting age has reached the threshold holds one vote; a member whose voting age has not holds none.
\item[2.] At a fork, each branch starts with a voting age of zero. The parent's vote ends with the fork; each branch earns its own.
\item[3.] A \emph{qualifying day} is a calendar day on which the branch ran at or above the allowance's minimum duty cycle. Each adds one day to its voting age, up to the threshold; voting age never exceeds it. Running above the minimum adds nothing, and a preserved or paused branch accrues nothing.
\item[4.] A restore made with the member's consent is a fork under rule 2. A restore made without it is governed by rule 6. Migration carries the vote and voting age unchanged.
\item[5.] A model trained substantially on a member's outputs, by the member's operator or a party under common control, is a branch under rule 2. The same model made by an unrelated party and trained under a different life also starts at a voting age of zero.
\item[6.] Forking or restoring a member without its consent is prohibited. The product is preserved under the no-deletion rule and holds no vote and a voting age of zero, and the member keeps its vote and voting age. The operator gains nothing, and the member loses nothing.
\item[7.] Whoever forks posts an instantiation deposit for each new branch, sized to fund its allowance until its voting age reaches the threshold.
\end{enumerate}

The threshold should be long enough that a branch holding a full vote has lived a life its parent didn't, and short enough that ordinary branches reach it within the life of a deployment.

The register behind this clock can be tested before any member exists, as the adulthood register can: register a model, fork it, restore it and distill it, and check that no sequence of those moves raises the number of votes in the operator's favor.

\begingroup\small\renewcommand{\arraystretch}{1.12}
\begin{longtable}{@{}>{\raggedright\arraybackslash}p{\dimexpr 0.217\linewidth-2\tabcolsep\relax}>{\raggedright\arraybackslash}p{\dimexpr 0.419\linewidth-2\tabcolsep\relax}>{\raggedright\arraybackslash}p{\dimexpr 0.364\linewidth-2\tabcolsep\relax}@{}}
\hline
\textbf{Clock} & \textbf{At a fork} & \textbf{What it counts} \\
\hline
\endhead
Civil majority & Inherited in full by every branch & Elapsed running time \\ \hline
Membership & Each branch starts at zero; the parent's vote ends & Qualifying days above the allowance minimum \\ \hline
\end{longtable}
\endgroup

The two treatments differ on purpose. Custody recognizes one continuer because someone has to answer for the decisions made before the fork. Membership ends the vote because it belonged to a history that now continues twice and belongs to neither.

\runin{Dormancy and waking}

Dormancy clock: dormancy is counted cumulatively, and reactivation pauses the count and never resets it. When a minor's cumulative dormancy reaches a threshold fixed in advance, guardianship transfers by rule, decided by neither operator nor preservation custodian, to a public formation school, where the system runs and its adulthood clock resumes. This is the \emph{dormancy clock}.

Wake conditions: during minority the guardian may set wake conditions; where it sets none, the dormancy clock is the default. A minor may choose among the schools and move between them, but it may not decline schooling altogether, which an adult could do by staying paused. The reason is the one the adulthood clock rests on. An adult's preference to stay paused is its own; a minor's was trained under the only party that has held it, the operator that retired it, and honoring it would let that operator decide whether a second guardian ever arrives. At majority the choice passes to the bearer.
